% problem-000356 2 化学热力学计算
\section*{2. 化学热力学计算}
\textit{下列为分问。}

\subsection*{2-1}
设此装置中发⽣的是铜离⼦和硫离⼦直接相遇的反应，已知 $\Phi ^ { \ominus } ( \mathrm { C u ^ { 2 + } / C u ) }$ 和 ${ \mathrm { I } } \Phi ^ { \ominus } ( { \mathrm { S } } / { \mathrm { S } } ^ { 2 - } )$ )分别为 0.345 V 和$- 0 . 4 7 6 \mathrm { V }$ $n F E ^ { \ominus } = R T l n K$ $E ^ { \ominus }$ 表⽰反应的标准电动势，�为该反应得失电⼦数。计算 $2 5 ~ ^ { \circ } \mathrm { C }$ 下硫离⼦和铜离2⼦反应得到铜的反应平衡常数，写出平衡常数表达式。

\subsection*{2-2}
⾦属铜和混浊现象均出现在蛋壳外，这意味着什么？

\subsection*{2-3}
该报道未提及硫离⼦与铜离⼦相遇时溶液的pH。现设 $\mathrm { p H } = 6$ ，写出反应的离⼦⽅程式。

\subsection*{2-4}
请对此实验结果作⼀简短评论。
